%% LyX 2.4.4 created this file.  For more info, see https://www.lyx.org/.
%% Do not edit unless you really know what you are doing.
\documentclass[english]{article}
\usepackage[T1]{fontenc}
\usepackage[latin9]{inputenc}
\usepackage{enumitem}
\usepackage{amsmath}
\usepackage{amsthm}
\usepackage{amssymb}
\usepackage{geometry}
\geometry{verbose,tmargin=1in,bmargin=1in,lmargin=1in,rmargin=1in}

\makeatletter
%%%%%%%%%%%%%%%%%%%%%%%%%%%%%% Textclass specific LaTeX commands.
\numberwithin{equation}{section}
\numberwithin{figure}{section}
\newlength{\lyxlabelwidth}      % auxiliary length 

%%%%%%%%%%%%%%%%%%%%%%%%%%%%%% User specified LaTeX commands.
\usepackage{fancyhdr}
\usepackage{lastpage}
\pagestyle{fancy}
\usepackage{enumitem}
\usepackage{ifthen}
\everymath=\expandafter{\the\everymath\displaystyle}
%\DeclareMathSizes{10}{10}{10}{10}
\usepackage{multicol}

\usepackage{tikz}
\usepackage{tikz-3dplot}
\usetikzlibrary{calc,arrows}

\usetikzlibrary{patterns}
\usetikzlibrary{plotmarks}
\usepackage{pgfplots}
\pgfplotsset{compat=newest}

\usepackage{calc}
\usepackage[nomessages]{fp}

\usepackage{datetime}
\usepackage[super]{nth}
% Added by lyx2lyx
\setlength{\parskip}{\smallskipamount}
\setlength{\parindent}{0pt}

\makeatother

\usepackage{babel}
\begin{document}
\renewcommand{\author}{Dr.\,James Ashe}

\newcommand{\classNum}{336}

\newcommand{\classSec}{A}

\newcommand{\examType}{HW 2.2 Solutions}

\renewcommand{\date}{\formatdate{4}{10}{2013}}

%%First page's header%%%%%%%%%%%%%%%%%%%%%%
\fancypagestyle{firststyle} %%header settings for first pagr
{
	\fancyhead[L]{MTH \classNum{}(\classSec{}) \author{}\\
	\textbf{\examType{}} ~\date{}}
	\fancyhead[C]{}
	\fancyhead[R]{\textbf{Name:\hspace{4cm}}}
	\setlength{\headsep}{14pt}
}
%%%%%%%%%%%%%%%%%%%%%%%%%%%%%%%%%%%%%%%%%%

%%Next page's header%%%%%%%%%%%%%%%%%%%%%%%
\setlist{nolistsep}
\lhead{\classNum{}(\classSec{})} 
\chead{\textbf{Name:}} 
\cfoot{\thepage{}~of~\pageref{LastPage}} 
\setlength{\headsep}{5pt} 
%%%%%%%%%%%%%%%%%%%%%%%%%%%%%%%%%%%%%%%%%%%

%%Various settings; see the preamble for other settings%%%%%
%\setenumerate[0]{leftmargin=12pt,itemindent=0pt,labelwidth=0pt}
\setlist{topsep=.5pt}
\renewcommand{\labelenumi}{\textbf{\arabic{enumi}.}}
\renewcommand{\labelenumii}{\textbf{\alph{enumii}.}}
\thispagestyle{firststyle}
%%%%%%%%%%%%%%%%%%%%%%%%%%%%%%%%%%%%%%%%%%

\newcommand{\xmax}{0}
\newcommand{\xmin}{-\xmax}
\newcommand{\xstep}{0}
\newcommand{\ymax}{0}
\newcommand{\ymin}{-\ymax}
\newcommand{\ystep}{0} 
\newcommand{\dspace}{\vspace{1cm}}
\newcommand{\xa}{0}
\newcommand{\xb}{0}
\newcommand{\ya}{1}
\newcommand{\yb}{1}

\small
\begin{enumerate}[resume]
\item Using the properties of linear transformations.

\begin{enumerate}[resume]
\item $T\left(2\begin{bmatrix}\begin{array}{r}
2\\
-1\\
1
\end{array}\end{bmatrix}\right)=2T\left(\begin{bmatrix}\begin{array}{r}
2\\
-1\\
1
\end{array}\end{bmatrix}\right)=2\begin{bmatrix}\begin{array}{r}
3\\
0
\end{array}\end{bmatrix}=\begin{bmatrix}\begin{array}{r}
6\\
0
\end{array}\end{bmatrix}$
\item $T\left(\begin{bmatrix}\begin{array}{r}
3\\
6\\
3
\end{array}\end{bmatrix}\right)=3T\left(\begin{bmatrix}\begin{array}{r}
1\\
2\\
1
\end{array}\end{bmatrix}\right)=3\begin{bmatrix}\begin{array}{r}
-1\\
2
\end{array}\end{bmatrix}=\begin{bmatrix}\begin{array}{r}
-3\\
6
\end{array}\end{bmatrix}$
\item $T\left(\begin{bmatrix}\begin{array}{r}
-1\\
3\\
0
\end{array}\end{bmatrix}\right)=T\left(\begin{bmatrix}\begin{array}{r}
1\\
2\\
1
\end{array}\end{bmatrix}\right)-T\left(\begin{bmatrix}\begin{array}{r}
2\\
-1\\
1
\end{array}\end{bmatrix}\right)=\begin{bmatrix}\begin{array}{r}
-1\\
2
\end{array}\end{bmatrix}-\begin{bmatrix}\begin{array}{r}
3\\
0
\end{array}\end{bmatrix}=\begin{bmatrix}\begin{array}{r}
-4\\
2
\end{array}\end{bmatrix}$
\end{enumerate}
\item Note that $x_{1}+2x_{2}+x_{3}=(1,2,1)\cdot(x_{1},x_{2},x_{3})$ and
$3x_{1}-x_{2}-x_{3}=(3,-1,-1)\cdot(x_{1},x_{2},x_{3})$. Also note
that $A\in\mathcal{M}_{2\times3}$ so that $A\mathbf{x}\in\mathcal{M}_{2\times1}$.
So 
\[
A=\begin{bmatrix}\begin{array}{rrr}
1 & 2 & 1\\
3 & -1 & -1
\end{array}\end{bmatrix}
\]
 Alternatively, $T\left(\begin{bmatrix}\begin{array}{r}
1\\
0\\
0
\end{array}\end{bmatrix}\right)=\begin{bmatrix}\begin{array}{r}
1\\
3
\end{array}\end{bmatrix}$, $T\left(\begin{bmatrix}\begin{array}{r}
0\\
1\\
0
\end{array}\end{bmatrix}\right)=\begin{bmatrix}\begin{array}{r}
2\\
-1
\end{array}\end{bmatrix}$, and $T\left(\begin{bmatrix}\begin{array}{r}
0\\
0\\
1
\end{array}\end{bmatrix}\right)=\begin{bmatrix}\begin{array}{r}
1\\
-1
\end{array}\end{bmatrix}$ gives the \nth{1}, \nth{2}, and \nth{3} column of $A$ respectively. 
\item \vspace{0cm}

\begin{enumerate}
\item The \nth{1} column is $T\left(\begin{bmatrix}\begin{array}{r}
1\\
0
\end{array}\end{bmatrix}\right)=\begin{bmatrix}\begin{array}{r}
2\\
-3
\end{array}\end{bmatrix}$. \\
For the \nth{2} column $T\left(\begin{bmatrix}\begin{array}{r}
0\\
1
\end{array}\end{bmatrix}\right)=T\left(\begin{bmatrix}\begin{array}{r}
2\\
1
\end{array}\end{bmatrix}\right)-T\left(\begin{bmatrix}\begin{array}{r}
2\\
0
\end{array}\end{bmatrix}\right)=\begin{bmatrix}\begin{array}{r}
-1\\
1
\end{array}\end{bmatrix}-\begin{bmatrix}\begin{array}{r}
4\\
-6
\end{array}\end{bmatrix}=\begin{bmatrix}\begin{array}{r}
-5\\
7
\end{array}\end{bmatrix}$. So 
\[
A=\begin{bmatrix}\begin{array}{rr}
2 & -5\\
-3 & 7
\end{array}\end{bmatrix}
\]
\end{enumerate}
\begin{enumerate}[resume]
\item The \nth{2} column is $T\left(\begin{bmatrix}\begin{array}{r}
0\\
1
\end{array}\end{bmatrix}\right)=\begin{bmatrix}\begin{array}{r}
1\\
-3
\end{array}\end{bmatrix}$. \\
The \nth{1} column is $T\left(\begin{bmatrix}\begin{array}{r}
1\\
0
\end{array}\end{bmatrix}\right)=\frac{1}{2}\left[T\left(\begin{bmatrix}\begin{array}{r}
2\\
1
\end{array}\end{bmatrix}\right)-T\left(\begin{bmatrix}\begin{array}{r}
0\\
1
\end{array}\end{bmatrix}\right)\right]=\frac{1}{2}\left[\begin{bmatrix}\begin{array}{r}
5\\
3
\end{array}\end{bmatrix}-\begin{bmatrix}\begin{array}{r}
1\\
-3
\end{array}\end{bmatrix}\right]=\frac{1}{2}\begin{bmatrix}\begin{array}{r}
4\\
6
\end{array}\end{bmatrix}=\begin{bmatrix}\begin{array}{r}
2\\
3
\end{array}\end{bmatrix}$. So 
\[
A=\begin{bmatrix}\begin{array}{rr}
2 & 1\\
3 & -3
\end{array}\end{bmatrix}
\]
\item $2T\left(\begin{bmatrix}\begin{array}{r}
1\\
0
\end{array}\end{bmatrix}\right)=T\left(\begin{bmatrix}\begin{array}{r}
2\\
0
\end{array}\end{bmatrix}\right)=T\left(\begin{bmatrix}\begin{array}{r}
1\\
1
\end{array}\end{bmatrix}\right)+T\left(\begin{bmatrix}\begin{array}{r}
1\\
-1
\end{array}\end{bmatrix}\right)=\begin{bmatrix}\begin{array}{r}
3\\
3
\end{array}\end{bmatrix}+\begin{bmatrix}\begin{array}{r}
-1\\
1
\end{array}\end{bmatrix}=\begin{bmatrix}\begin{array}{r}
2\\
4
\end{array}\end{bmatrix}$; so the \nth{1} column is $\begin{bmatrix}\begin{array}{r}
1\\
2
\end{array}\end{bmatrix}$ (divide by $2$). \\
$2T\left(\begin{bmatrix}\begin{array}{r}
0\\
1
\end{array}\end{bmatrix}\right)=T\left(\begin{bmatrix}\begin{array}{r}
1\\
1
\end{array}\end{bmatrix}\right)-T\left(\begin{bmatrix}\begin{array}{r}
1\\
-1
\end{array}\end{bmatrix}\right)=\begin{bmatrix}\begin{array}{r}
3\\
3
\end{array}\end{bmatrix}-\begin{bmatrix}\begin{array}{r}
-1\\
1
\end{array}\end{bmatrix}=\begin{bmatrix}\begin{array}{r}
4\\
2
\end{array}\end{bmatrix}$; so the \nth{2} column is $\begin{bmatrix}\begin{array}{r}
2\\
1
\end{array}\end{bmatrix}$. Thus
\[
A=\begin{bmatrix}\begin{array}{rr}
1 & 2\\
2 & 1
\end{array}\end{bmatrix}
\]
\end{enumerate}
\item Recall that to prove that $T$ is a linear transformation it suffices
to find a standard matrix, $A\in\mathcal{M}_{m\times n}$ for $\mu_{A}:\mathbb{R}^{n}\rightarrow\mathbb{R}^{m}$.
To prove that it is not a standard matrix you need to show that it
fails at least one of the properties.

\begin{enumerate}
\item $T$ is not a linear transformation.

\begin{proof}
$-1\cdot T\left(\begin{bmatrix}\begin{array}{r}
0\\
1
\end{array}\end{bmatrix}\right)=-\begin{bmatrix}\begin{array}{r}
2\\
1
\end{array}\end{bmatrix}=\begin{bmatrix}\begin{array}{r}
-2\\
-1
\end{array}\end{bmatrix}\neq T\left(\begin{bmatrix}\begin{array}{r}
0\\
-1
\end{array}\end{bmatrix}\right)=\begin{bmatrix}\begin{array}{r}
-2\\
1
\end{array}\end{bmatrix}$.
\end{proof}
\end{enumerate}
\begin{enumerate}[resume]
\item $T$ is a linear transformation. 

\begin{proof}
$T\left(\begin{bmatrix}\begin{array}{r}
1\\
0
\end{array}\end{bmatrix}\right)=\begin{bmatrix}\begin{array}{r}
1\\
0
\end{array}\end{bmatrix}$ and $T\left(\begin{bmatrix}\begin{array}{r}
0\\
1
\end{array}\end{bmatrix}\right)=\begin{bmatrix}\begin{array}{r}
2\\
0
\end{array}\end{bmatrix}$ gives the \nth{1} and \nth{2} column respectively of $A=\begin{bmatrix}\begin{array}{rr}
1 & 2\\
0 & 0
\end{array}\end{bmatrix}$, and $\begin{bmatrix}\begin{array}{rr}
1 & 2\\
0 & 0
\end{array}\end{bmatrix}\begin{bmatrix}\begin{array}{r}
x_{1}\\
x_{2}
\end{array}\end{bmatrix}=\begin{bmatrix}\begin{array}{r}
x_{1}+2x_{2}\\
0
\end{array}\end{bmatrix}$ so $\mu_{\mbox{A}}=T$ which shows that $T$ is linear. 
\end{proof}
\end{enumerate}
\item [\textbf{6.}]

\begin{enumerate}
\item $T\left(\begin{bmatrix}\begin{array}{r}
1\\
0
\end{array}\end{bmatrix}\right)=\begin{bmatrix}\begin{array}{r}
0\\
1
\end{array}\end{bmatrix}$ and $T\left(\begin{bmatrix}\begin{array}{r}
0\\
1
\end{array}\end{bmatrix}\right)=\begin{bmatrix}\begin{array}{r}
-1\\
0
\end{array}\end{bmatrix}$$\Rightarrow A_{T}=\begin{bmatrix}\begin{array}{rr}
0 & -1\\
1 & 0
\end{array}\end{bmatrix}$\\
It is a reflection across $x_{1}+x_{2}=0\Rightarrow x_{1}=-x_{2}$.
So then $S\left(\begin{bmatrix}\begin{array}{r}
1\\
0
\end{array}\end{bmatrix}\right)=\begin{bmatrix}\begin{array}{r}
0\\
-1
\end{array}\end{bmatrix}$ and $S\left(\begin{bmatrix}\begin{array}{r}
0\\
1
\end{array}\end{bmatrix}\right)=\begin{bmatrix}\begin{array}{r}
-1\\
0
\end{array}\end{bmatrix}$$\Rightarrow A_{s}=\begin{bmatrix}\begin{array}{rr}
0 & -1\\
-1 & 0
\end{array}\end{bmatrix}$.
\end{enumerate}
\begin{enumerate}[resume]
\item $T\circ S$ notates the composition function, i.e., the points are
mapped first by $S$ and then by $T$. So $\left(T\circ S\right)\left(\begin{bmatrix}\begin{array}{r}
1\\
0
\end{array}\end{bmatrix}\right)=T\left(\begin{bmatrix}\begin{array}{r}
0\\
-1
\end{array}\end{bmatrix}\right)=\begin{bmatrix}\begin{array}{rr}
0 & -1\\
1 & 0
\end{array}\end{bmatrix}\begin{bmatrix}\begin{array}{r}
0\\
-1
\end{array}\end{bmatrix}=\begin{bmatrix}1\\
0
\end{bmatrix}$ and $\left(T\circ S\right)\left(\begin{bmatrix}\begin{array}{r}
0\\
1
\end{array}\end{bmatrix}\right)=T\left(\begin{bmatrix}\begin{array}{r}
-1\\
0
\end{array}\end{bmatrix}\right)=\begin{bmatrix}\begin{array}{rr}
0 & -1\\
1 & 0
\end{array}\end{bmatrix}\begin{bmatrix}\begin{array}{r}
-1\\
0
\end{array}\end{bmatrix}=\begin{bmatrix}0\\
-1
\end{bmatrix}$. So the standard matrix is $A_{T\circ S}=\begin{bmatrix}\begin{array}{rr}
1 & 0\\
0 & -1
\end{array}\end{bmatrix}$. Note that $TS=A_{T\circ S}$.
\item $\left(S\circ T\right)\left(\begin{bmatrix}\begin{array}{r}
1\\
0
\end{array}\end{bmatrix}\right)=T\left(\begin{bmatrix}\begin{array}{r}
0\\
1
\end{array}\end{bmatrix}\right)=\begin{bmatrix}\begin{array}{rr}
0 & -1\\
-1 & 0
\end{array}\end{bmatrix}\begin{bmatrix}\begin{array}{r}
0\\
1
\end{array}\end{bmatrix}=\begin{bmatrix}-1\\
0
\end{bmatrix}$ and $\left(S\circ T\right)\left(\begin{bmatrix}\begin{array}{r}
0\\
1
\end{array}\end{bmatrix}\right)=T\left(\begin{bmatrix}\begin{array}{r}
-1\\
0
\end{array}\end{bmatrix}\right)=\begin{bmatrix}\begin{array}{rr}
0 & -1\\
-1 & 0
\end{array}\end{bmatrix}\begin{bmatrix}\begin{array}{r}
-1\\
0
\end{array}\end{bmatrix}=\begin{bmatrix}0\\
1
\end{bmatrix}$. So the standard matrix is $A_{T\circ S}=\begin{bmatrix}\begin{array}{rr}
-1 & 0\\
0 & 1
\end{array}\end{bmatrix}$. Note that $ST=A_{S\circ T}$.
\end{enumerate}
\end{enumerate}

\end{document}
